% ============================================================================
%  Quantex 内部讲义 · 《量化金融工程》统一导言区
%  ---------------------------------------------------------------------------
%  引擎：XeLaTeX          文档类：ctexbook（在 main.tex 中声明）
%  本文件被 main.tex 与各章独立编译脚本（tools/check_chapter.sh）共同引用，
%  所有排版样式、数学宏、彩色盒子、算法与代码样式均集中于此，便于统一维护。
%  修改样式只需改这里一处，全书与单章 PDF 同时生效。
% ============================================================================

% ---------------------------------------------------------------- 引擎与数学
\usepackage{iftex}
\RequireXeTeX
\usepackage{amsmath,amssymb,amsthm,mathtools}
\usepackage{bm}
\usepackage{mathrsfs}

% ---------------------------------------------------------------- 颜色（须先于 tikz/pgfplots 载入，避免选项冲突）
\usepackage[table]{xcolor}

% ---------------------------------------------------------------- 版面与排版
\usepackage[
  a4paper, left=2.5cm, right=2.5cm, top=2.6cm, bottom=2.4cm,
  headheight=16pt, headsep=12pt, footskip=1.1cm
]{geometry}
\usepackage{graphicx}
\usepackage{array,booktabs,multirow,makecell,tabularx,longtable}
\usepackage{enumitem}
\usepackage{caption}
\usepackage{tocloft}
\usepackage{fancyhdr}
\usepackage{float}
\usepackage{algorithm}
\usepackage{algpseudocode}
\usepackage{hyperref}   % 先载入，稍后由 \hypersetup 统一配置

% ---------------------------------------------------------------- 图形与绘图
\usepackage{tikz}
\usetikzlibrary{
  arrows.meta, positioning, calc, patterns,
  decorations.pathreplacing, decorations.markings,
  shapes.geometric, shapes.misc, fit, backgrounds, matrix
}
\usepackage{pgfplots}
\pgfplotsset{compat=1.18}
\usepgfplotslibrary{fillbetween}

% ---------------------------------------------------------------- 颜色与盒子
\usepackage{fontawesome5}
\usepackage{tcolorbox}
\tcbuselibrary{skins,breakable,theorems,hooks}
\usepackage{listings}

% ============================================================================
%  配色方案
% ============================================================================
\definecolor{quantblue}{RGB}{30,64,175}     % 主色：量化蓝
\definecolor{quantnavy}{RGB}{15,32,84}      % 深蓝：封面底色
\definecolor{quantgreen}{RGB}{5,150,105}    % 辅色：收益/正向
\definecolor{quantred}{RGB}{185,28,28}      % 警示色：风险/负向
\definecolor{quantgold}{RGB}{176,122,30}    % 点缀色：金色分隔线
\definecolor{quantpurple}{RGB}{109,40,217}  % 第四色：概率/测度
\definecolor{inkgray}{RGB}{55,65,81}        % 正文灰：标题层级
\definecolor{codebg}{RGB}{247,248,250}      % 代码底色
\definecolor{boxbg}{RGB}{239,246,255}       % 要点底色
\definecolor{warnbg}{RGB}{254,242,242}      % 警示底色
\definecolor{appbg}{RGB}{240,253,244}       % 应用底色
\definecolor{intubg}{RGB}{253,250,235}      % 直觉底色

% ============================================================================
%  中文字体（macOS 字体回退：宋体正文 / 黑体标题 / 楷体强调）
% ============================================================================
\setCJKmainfont{Songti SC}[BoldFont=Songti SC Bold, ItalicFont=Kaiti SC]
\setCJKsansfont{Heiti SC}[BoldFont=Heiti SC Medium]
\setCJKmonofont{Heiti SC}
\setlength{\parindent}{0pt}
\setlength{\parskip}{0.45em}
\linespread{1.14}

% ============================================================================
%  章节标题样式（ctexset）
%  注意：不覆盖 chapter/name，以便 \appendix 自动切换为「附录 A」
% ============================================================================
\ctexset{
  chapter = {
    format     = \raggedright\sffamily\bfseries\color{quantblue}%
                 \fontsize{25}{29}\selectfont,
    nameformat = \color{quantblue!45},
    aftername  = \hspace{0.35em},
    beforeskip = {0pt},
    afterskip  = {16pt},
    aftertitle = {\par\vspace{2pt}%
                  {\color{quantgold}\rule{\textwidth}{1.6pt}}\par\vspace{-2pt}},
    pagestyle  = fancy
  },
  section = {
    format     = \Large\sffamily\bfseries\color{quantblue},
    aftername  = \hspace{0.3em},
    beforeskip = {15pt plus 3pt minus 2pt},
    afterskip  = {7pt plus 1pt minus 1pt}
  },
  subsection = {
    format     = \large\sffamily\bfseries\color{inkgray},
    beforeskip = {11pt plus 2pt minus 2pt},
    afterskip  = {5pt plus 1pt minus 1pt}
  },
  subsubsection = {
    format     = \normalsize\sffamily\bfseries\color{inkgray!85},
    beforeskip = {9pt plus 2pt minus 2pt},
    afterskip  = {4pt plus 1pt minus 1pt}
  }
}

% ============================================================================
%  目录样式
% ============================================================================
\renewcommand{\cftchapfont}{\sffamily\bfseries\color{quantblue}}
\renewcommand{\cftchappagefont}{\sffamily\bfseries\color{quantblue}}
\renewcommand{\cftsecfont}{\sffamily\color{inkgray}}
\renewcommand{\cftsecpagefont}{\sffamily\color{inkgray}}
\renewcommand{\cftsubsecfont}{\sffamily\small\color{inkgray!80}}
\renewcommand{\cftsubsecpagefont}{\sffamily\small\color{inkgray!80}}
\renewcommand{\cftchapleader}{\cftdotfill{\cftdotsep}}
\setlength{\cftbeforechapskip}{7pt}
\setlength{\cftbeforesecskip}{2pt}
\setlength{\cftbeforesubsecskip}{1pt}
\setcounter{tocdepth}{2}
\setcounter{secnumdepth}{3}

% 目录标题「目录」本身
\renewcommand{\contentsname}{目\hspace{0.5em}录}

% ============================================================================
%  页眉页脚
% ============================================================================
\pagestyle{fancy}
\fancyhf{}
\fancyhead[LE,RO]{\sffamily\bfseries\color{quantblue}\thepage}
\fancyhead[RE]{\sffamily\small\color{inkgray}\nouppercase{\leftmark}}
\fancyhead[LO]{\sffamily\small\color{inkgray}\nouppercase{\rightmark}}
\fancyfoot[C]{\sffamily\scriptsize\color{inkgray!65}%
  Quantex 内部讲义 · 量化金融工程 · Quantitative Financial Engineering}
\renewcommand{\headrulewidth}{0pt}
\renewcommand{\footrulewidth}{0pt}
\renewcommand{\headrule}{{\color{quantblue!22}\hrule width\headwidth height 0.9pt}}

% 前置部分（前言/目录）使用简洁页脚
\fancypagestyle{frontstyle}{%
  \fancyhf{}%
  \fancyfoot[C]{\sffamily\small\color{inkgray!70}\thepage}%
  \renewcommand{\headrule}{}%
}

% ============================================================================
%  超链接
% ============================================================================
\hypersetup{
  colorlinks   = true,
  linkcolor    = quantblue,
  citecolor    = quantgreen,
  urlcolor     = quantblue!85!black,
  linktoc      = all,
  bookmarksnumbered = true,
  bookmarksopen     = false,
  unicode      = true,
  pdfencoding  = auto,
  pdftitle     = {量化金融工程 · Quantitative Financial Engineering},
  pdfauthor    = {Quantex},
  pdfsubject   = {Quantex 内部讲义 · 量化金融 · 衍生品定价 · 因子投资 · 风险管理},
  pdfkeywords  = {量化金融, Q Quant, P Quant, 衍生品定价, 随机过程, 因子模型, 组合优化}
}
% 标题中如不得已出现数学，保证 PDF 书签可读
\pdfstringdefDisableCommands{%
  \def\ensuremath#1{#1}%
  \def\frac#1#2{#1/#2}%
  \def\text#1{#1}%
  \def\mathrm#1{#1}%
  \def\mathbb#1{#1}%
}

% ============================================================================
%  交叉引用（cleveref，中文名称）
% ============================================================================
\usepackage{cleveref}
\crefname{chapter}{章}{章}
\crefname{section}{节}{节}
\crefname{subsection}{小节}{小节}
\crefname{equation}{式}{式}
\crefname{figure}{图}{图}
\crefname{table}{表}{表}
\crefname{lstlisting}{代码}{代码}
\crefname{algorithm}{算法}{算法}
\crefname{definition}{定义}{定义}
\crefname{theorem}{定理}{定理}
\crefname{proposition}{命题}{命题}
\crefname{lemma}{引理}{引理}
\crefname{corollary}{推论}{推论}
\crefname{example}{例}{例}
\crefname{remark}{注记}{注记}
\crefname{assumption}{假设}{假设}
\crefformat{equation}{式~(#2#1#3)}
\crefformat{lstlisting}{代码~#2#1#3}
\crefrangeformat{equation}{式~(#3#1#4)---(#5#2#6)}

% ============================================================================
%  定理类环境
%  · definition / example 使用彩色盒子（视觉上更醒目，编号按章）
%  · theorem / proposition / lemma / corollary / remark / assumption 使用
%    自定义定理样式（较轻量，避免整页都是色块）
% ============================================================================
\newtheoremstyle{quant}
  {9pt}          % 上方间距
  {9pt}          % 下方间距
  {\normalfont}  % 正文
  {0pt}          % 缩进
  {\bfseries\sffamily\color{quantblue}}  % 标题
  {}             % 标题后标点
  {0.6em}        % 标题后间距
  {}
\theoremstyle{quant}
\newtheorem{theorem}{定理}[chapter]
\newtheorem{proposition}{命题}[chapter]
\newtheorem{lemma}{引理}[chapter]
\newtheorem{corollary}{推论}[chapter]
\newtheorem{remark}{注记}[chapter]
\newtheorem{assumption}{假设}[chapter]
\newtheorem{exercise}{思考题}[chapter]

\renewcommand{\proofname}{\sffamily\bfseries 证明}
\renewcommand{\qedsymbol}{$\blacksquare$}

% --- 定义盒 ---------------------------------------------------------------
\newtcolorbox[auto counter, number within=chapter]{definition}[1][]{
  enhanced, breakable,
  colback=boxbg, colframe=quantblue!55, boxrule=0.6pt, arc=1.5mm,
  left=3mm, right=3mm, top=2mm, bottom=2mm,
  fonttitle=\bfseries\sffamily,
  title={定义~\thetcbcounter\ \ \textnormal{#1}},
  attach boxed title to top left={yshift=-2mm, xshift=3mm},
  boxed title style={colback=quantblue!55, size=small, arc=1mm, boxrule=0pt},
  coltitle=black
}
\crefname{tcb@cnt@definition}{定义}{定义}
\Crefname{tcb@cnt@definition}{定义}{定义}

% --- 例盒 -----------------------------------------------------------------
\newtcolorbox[auto counter, number within=chapter]{example}[1][]{
  enhanced, breakable,
  colback=quantgreen!5, colframe=quantgreen!55, boxrule=0.6pt, arc=1.5mm,
  left=3mm, right=3mm, top=2mm, bottom=2mm,
  fonttitle=\bfseries\sffamily,
  title={例~\thetcbcounter\ \ \textnormal{#1}},
  attach boxed title to top left={yshift=-2mm, xshift=3mm},
  boxed title style={colback=quantgreen!55, size=small, arc=1mm, boxrule=0pt},
  coltitle=black
}
\crefname{tcb@cnt@example}{例}{例}
\Crefname{tcb@cnt@example}{例}{例}

% ============================================================================
%  彩色提示盒
% ============================================================================
% 要点盒（按章编号，可用 \label + \cref 引用）
\newtcolorbox[auto counter, number within=chapter]{keybox}[1][]{
  enhanced, breakable,
  colback=boxbg, colframe=quantblue!70, boxrule=0.8pt, arc=1.5mm,
  left=3mm, right=3mm, top=2mm, bottom=2mm,
  fonttitle=\bfseries\sffamily\color{white},
  coltitle=white,
  title={\faLightbulb\hspace{0.4em}核心要点~\thetcbcounter\ \ #1},
  attach boxed title to top left={yshift=-2mm, xshift=3mm},
  boxed title style={colback=quantblue!80, size=small, arc=1mm, boxrule=0pt}
}
\crefname{tcb@cnt@keybox}{要点}{要点}
\Crefname{tcb@cnt@keybox}{要点}{要点}

% 风险警示盒
\newtcolorbox{warningbox}[1][]{
  enhanced, breakable,
  colback=warnbg, colframe=quantred!70, boxrule=0.8pt, arc=1.5mm,
  left=3mm, right=3mm, top=2mm, bottom=2mm,
  fonttitle=\bfseries\sffamily\color{white}, coltitle=white,
  title={\faExclamationTriangle\hspace{0.4em}风险警示\ \ #1},
  attach boxed title to top left={yshift=-2mm, xshift=3mm},
  boxed title style={colback=quantred!80, size=small, arc=1mm, boxrule=0pt}
}

% 实务应用盒
\newtcolorbox{appbox}[1][]{
  enhanced, breakable,
  colback=appbg, colframe=quantgreen!65, boxrule=0.8pt, arc=1.5mm,
  left=3mm, right=3mm, top=2mm, bottom=2mm,
  fonttitle=\bfseries\sffamily\color{white}, coltitle=white,
  title={\faChartLine\hspace{0.4em}实务应用\ \ #1},
  attach boxed title to top left={yshift=-2mm, xshift=3mm},
  boxed title style={colback=quantgreen!70, size=small, arc=1mm, boxrule=0pt}
}

% 直觉理解盒
\newtcolorbox{intuitionbox}[1][]{
  enhanced, breakable,
  colback=intubg, colframe=quantgold!65, boxrule=0.8pt, arc=1.5mm,
  left=3mm, right=3mm, top=2mm, bottom=2mm,
  fonttitle=\bfseries\sffamily, coltitle=black,
  title={\faBrain\hspace{0.4em}直觉理解\ \ #1},
  attach boxed title to top left={yshift=-2mm, xshift=3mm},
  boxed title style={colback=quantgold!45, size=small, arc=1mm, boxrule=0pt}
}

% 补充说明盒
\newtcolorbox{notebox}[1][]{
  enhanced, breakable,
  colback=white, colframe=inkgray!45, boxrule=0.6pt, arc=1.5mm,
  left=3mm, right=3mm, top=2mm, bottom=2mm,
  fonttitle=\bfseries\sffamily, coltitle=inkgray,
  title={\faInfoCircle\hspace{0.4em}补充说明\ \ #1},
  attach boxed title to top left={yshift=-2mm, xshift=3mm},
  boxed title style={colback=inkgray!12, size=small, arc=1mm, boxrule=0pt}
}

% 考点盒（用于资格证章节）
\newtcolorbox{exambox}[1][]{
  enhanced, breakable,
  colback=quantpurple!5, colframe=quantpurple!60, boxrule=0.8pt, arc=1.5mm,
  left=3mm, right=3mm, top=2mm, bottom=2mm,
  fonttitle=\bfseries\sffamily\color{white}, coltitle=white,
  title={\faGraduationCap\hspace{0.4em}考试要点\ \ #1},
  attach boxed title to top left={yshift=-2mm, xshift=3mm},
  boxed title style={colback=quantpurple!70, size=small, arc=1mm, boxrule=0pt}
}

% ============================================================================
%  章节收尾环境：本章小结 / 练习与思考 / 延伸阅读
% ============================================================================
\newtcolorbox{summarybox}[1][本章小结]{
  enhanced, breakable,
  colback=quantblue!4, colframe=quantblue!45, boxrule=0.8pt, arc=2mm,
  left=3.5mm, right=3.5mm, top=2.5mm, bottom=2.5mm,
  fonttitle=\bfseries\sffamily\large\color{quantblue},
  coltitle=quantblue, colbacktitle=quantblue!8,
  title={\faBookOpen\hspace{0.4em}#1},
  attach boxed title to top left={yshift=-2.5mm, xshift=4mm},
  boxed title style={colback=quantblue!8, size=small, arc=1mm, boxrule=0pt}
}

\newenvironment{exercises}[1][练习与思考]{%
  \par\vspace{6pt}%
  {\sffamily\bfseries\large\color{quantblue}\faPen\hspace{0.4em}#1}\par
  \vspace{4pt}%
  \begin{enumerate}[label=\textbf{\arabic*.}, leftmargin=1.9em,
                    itemsep=3pt, topsep=2pt, parsep=0pt]%
}{%
  \end{enumerate}\par\vspace{4pt}%
}

\newenvironment{furtherreading}[1][延伸阅读]{%
  \par\vspace{6pt}%
  {\sffamily\bfseries\large\color{quantblue}\faBook\hspace{0.4em}#1}\par
  \vspace{4pt}%
  \begin{itemize}[label=\textcolor{quantgold}{\small\faAngleRight},
                  leftmargin=1.5em, itemsep=2pt, topsep=2pt, parsep=0pt]%
}{%
  \end{itemize}\par\vspace{4pt}%
}

% 用于「延伸阅读」中的文献条目：作者《标题》，期刊/出版社，年份
\newcommand{\refitem}[1]{\item #1}

% ============================================================================
%  表格与图注样式
% ============================================================================
\captionsetup{
  font={small,sf},
  labelfont={bf,color=quantblue},
  labelsep=quad,
  skip=6pt,
  justification=centering
}
\captionsetup[table]{skip=4pt}
\captionsetup[algorithm]{font={small,sf},labelfont={bf,color=quantpurple},labelsep=quad}

% 列类型：定宽左/中/右对齐，以及 tabularx 弹性列
\newcolumntype{L}[1]{>{\raggedright\arraybackslash}p{#1}}
\newcolumntype{C}[1]{>{\centering\arraybackslash}p{#1}}
\newcolumntype{R}[1]{>{\raggedleft\arraybackslash}p{#1}}
\newcolumntype{Y}{>{\raggedright\arraybackslash}X}

\newcommand{\tablefont}{\small\setlength{\tabcolsep}{5.5pt}%
  \renewcommand{\arraystretch}{1.16}}

% 行内术语首次出现：中文（English）
\newcommand{\term}[1]{\textbf{#1}}

% ============================================================================
%  算法与代码样式
% ============================================================================
\floatname{algorithm}{算法}
\renewcommand{\algorithmicrequire}{\textbf{输入:}}
\renewcommand{\algorithmicensure}{\textbf{输出:}}
\renewcommand{\algorithmiccomment}[1]{\hfill\textcolor{quantgreen}{\itshape\small // #1}}
\renewcommand{\listalgorithmname}{算法目录}
\numberwithin{algorithm}{chapter}

\lstdefinestyle{pythonstyle}{
  language=Python,
  backgroundcolor=\color{codebg},
  basicstyle=\ttfamily\footnotesize,
  keywordstyle=\color{quantblue}\bfseries,
  commentstyle=\color{quantgreen}\itshape,
  stringstyle=\color{quantred},
  numbers=left, numberstyle=\tiny\color{inkgray!60}, numbersep=6pt,
  frame=single, framerule=0.4pt, rulecolor=\color{inkgray!35},
  breaklines=true, breakatwhitespace=false,
  showstringspaces=false, tabsize=2,
  captionpos=t, abovecaptionskip=4pt, belowcaptionskip=4pt,
  xleftmargin=14pt, framexleftmargin=14pt
}
\lstdefinestyle{soliditystyle}{
  language=Java,
  backgroundcolor=\color{codebg},
  basicstyle=\ttfamily\footnotesize,
  keywordstyle=\color{quantpurple}\bfseries,
  commentstyle=\color{quantgreen}\itshape,
  stringstyle=\color{quantred},
  numbers=left, numberstyle=\tiny\color{inkgray!60}, numbersep=6pt,
  frame=single, framerule=0.4pt, rulecolor=\color{inkgray!35},
  breaklines=true, showstringspaces=false, tabsize=4,
  captionpos=t, abovecaptionskip=4pt, belowcaptionskip=4pt,
  xleftmargin=14pt, framexleftmargin=14pt
}
\lstdefinestyle{shellstyle}{
  language=bash,
  backgroundcolor=\color{codebg},
  basicstyle=\ttfamily\footnotesize,
  keywordstyle=\color{quantblue}\bfseries,
  commentstyle=\color{quantgreen}\itshape,
  numbers=none,
  frame=single, framerule=0.4pt, rulecolor=\color{inkgray!35},
  breaklines=true, showstringspaces=false,
  captionpos=t, xleftmargin=14pt, framexleftmargin=14pt
}
\lstdefinestyle{plaincode}{
  backgroundcolor=\color{codebg},
  basicstyle=\ttfamily\footnotesize,
  numbers=left, numberstyle=\tiny\color{inkgray!60}, numbersep=6pt,
  frame=single, framerule=0.4pt, rulecolor=\color{inkgray!35},
  breaklines=true, showstringspaces=false,
  captionpos=t, xleftmargin=14pt, framexleftmargin=14pt
}

% ============================================================================
%  数学宏（全书统一，鼓励使用以保持符号一致）
% ============================================================================
\newcommand{\E}{\mathbb{E}}                        % 期望
\newcommand{\Prob}{\mathbb{P}}                     % 真实测度
\newcommand{\Qm}{\mathbb{Q}}                       % 风险中性测度
\newcommand{\Fm}{\mathcal{F}}                      % 信息流
\newcommand{\N}{\mathcal{N}}                       % 正态/鞅
\newcommand{\R}{\mathbb{R}}                        % 实数集
\newcommand{\Var}{\operatorname{Var}}              % 方差
\newcommand{\Cov}{\operatorname{Cov}}              % 协方差
\newcommand{\Corr}{\operatorname{Corr}}            % 相关系数
\newcommand{\sd}{\operatorname{sd}}                % 标准差
\newcommand{\se}{\operatorname{SE}}                % 标准误
\newcommand{\VaR}{\operatorname{VaR}}              % 在险价值
\newcommand{\ES}{\operatorname{ES}}                % 预期亏空
\newcommand{\abs}[1]{\left\lvert #1 \right\rvert}
\newcommand{\norm}[1]{\left\lVert #1 \right\rVert}
\newcommand{\inner}[2]{\left\langle #1,\, #2 \right\rangle}
\newcommand{\dd}{\mathrm{d}}                       % 微分 d
\newcommand{\ee}{\mathrm{e}}                       % 自然指数 e
\newcommand{\ii}{\mathrm{i}}                       % 虚数单位
\newcommand{\ind}[1]{\mathbf{1}_{\{#1\}}}          % 示性函数
\newcommand{\tr}{\operatorname{tr}}
\newcommand{\diag}{\operatorname{diag}}
\newcommand{\rank}{\operatorname{rank}}
\newcommand{\vect}[1]{\bm{#1}}                     % 向量/矩阵粗体
\newcommand{\pd}[2]{\frac{\partial #1}{\partial #2}}
\newcommand{\pdd}[2]{\frac{\partial^2 #1}{\partial #2^2}}
\DeclareMathOperator*{\argmax}{arg\,max}
\DeclareMathOperator*{\argmin}{arg\,min}
\DeclareMathOperator{\sign}{sgn}
\DeclareMathOperator{\law}{\mathcal{L}}
\DeclareMathOperator{\supp}{supp}

% 常用简写：收益率、波动率、相关系数、权重
\newcommand{\ret}{r}
\newcommand{\vol}{\sigma}
\newcommand{\wvec}{\vect{w}}
\newcommand{\muvec}{\bm{\mu}}
\newcommand{\Sigmamat}{\bm{\Sigma}}
\newcommand{\volvec}{\bm{\sigma}}                  % 波动率向量
\newcommand{\corrmat}{\bm{R}}                      % 相关矩阵
\newcommand{\ones}{\mathbf{1}}                     % 全 1 向量
\newcommand{\CVaR}{\operatorname{CVaR}}             % 条件在险价值（= ES）
\newcommand{\SR}{\operatorname{SR}}                 % 夏普比率
\newcommand{\IC}{\operatorname{IC}}                 % 信息系数
\newcommand{\IR}{\operatorname{IR}}                 % 信息比率
\newcommand{\Gm}{\mathcal{G}}                      % 辅助信息流
\newcommand{\Hm}{\mathcal{H}}                      % 辅助信息流
